\clearpage
\appendix

\section{预览}
\begin{studygoal}
  \begin{itemize}
    \item 能够用日语完成简单的自我介绍。
    \item 掌握判断句「A は B です」及疑问句「A は B ですか」。
    \item 记住初次见面时常用的寒暄表达。
  \end{itemize}
\end{studygoal}

\subsection*{核心词汇}

\begin{vocab}{はじめまして}
  \vocabentry{私}{わたし}{代词}{我}
  \vocabentry{学生}{がくせい}{名词}{学生}
  \vocabentry{先生}{せんせい}{名词}{老师；也用于称呼医生、律师等}
  \vocabentry{中国人}{ちゅうごくじん}{名词}{中国人}
  \vocabentry{日本語}{にほんご}{名词}{日语}
  \vocabentry{よろしく}{よろしく}{副词}{请多关照；请多指教}
\end{vocab}

\begin{vocabtable}{}
  大学 & だいがく & 名词 & 大学 & 大学の学生\\
  専攻 & せんこう & 名词 & 专业；专攻 & 専攻は数学です\\
  勉強する & べんきょうする & 动词 & 学习 & 日本語を勉強する\\
  友達 & ともだち & 名词 & 朋友 & 日本人の友達\\
\end{vocabtable}

\subsection*{假名与发音}

日语词汇第一次出现时，可以用
\furigana{漢字}{かんじ}
标注读音。连续句子不必为每个词都加注，通常只给生词或易错读音加注即可。

\begin{center}
  \renewcommand{\arraystretch}{1.5}
  \begin{tabular}{c|ccccc}
    \rowcolor{JPIndigo!16}
       & あ段 & い段 & う段 & え段 & お段\\
    \hline
    あ行 & \jp{あ} & \jp{い} & \jp{う} & \jp{え} & \jp{お}\\
    か行 & \jp{か} & \jp{き} & \jp{く} & \jp{け} & \jp{こ}\\
    さ行 & \jp{さ} & \jp{し} & \jp{す} & \jp{せ} & \jp{そ}\\
    た行 & \jp{た} & \jp{ち} & \jp{つ} & \jp{て} & \jp{と}\\
    な行 & \jp{な} & \jp{に} & \jp{ぬ} & \jp{ね} & \jp{の}\\
  \end{tabular}
\end{center}

\begin{pitfall}[发音提醒]
  \jp{「は」}作主题助词时读作 \jp{「わ」}，作普通假名时仍读作
  \jp{「は」}。例如：\jp{私は学生です。}
\end{pitfall}

\subsection*{基础语法}

\begin{grammar}{A は B です}
  这是日语中最基础的判断句，可以理解为“A 是 B”。其中
  \term{は}提示句子的主题，\term{です}使表达礼貌、完整。

  \[
    \text{A（主题）}+\text{\jp{は}}+\text{B（说明）}+\text{\jp{です}}
  \]

  否定形式为 \jp{「A は B ではありません」}，口语中也常说
  \jp{「A は B じゃありません」}。
\end{grammar}

\begin{example}
  \sentence{私は学生です。}{わたし は がくせい です。}{我是学生。}
  \tcblower
  \sentence{田中さんは先生ではありません。}
    {たなかさん は せんせい では ありません。}
    {田中先生不是老师。}
\end{example}

\begin{grammar}{A は B ですか}
  在礼貌判断句末尾加疑问助词 \term{か}，即可构成一般疑问句。
  日语书写中即使不写问号，也能通过句末的 \jp{「か」}表示疑问。
\end{grammar}

\begin{example}
  \sentence{李さんは大学生ですか。}
    {りさん は だいがくせい ですか。}
    {小李是大学生吗？}
  \tcblower
  \sentence{はい、大学生です。}{はい、だいがくせい です。}{是的，我是大学生。}
\end{example}

\subsection*{会话}

\begin{dialogue}
  \dialogueline{A}{はじめまして。田中です。}{初次见面，我是田中。}
  \dialogueline{B}{はじめまして。李です。}{初次见面，我姓李。}
  \dialogueline{A}{李さんは学生ですか。}{小李是学生吗？}
  \dialogueline{B}{はい、大学生です。よろしくお願いします。}
    {是的，我是大学生。请多关照。}
  \dialogueline{A}{こちらこそ、よろしくお願いします。}
    {我才要请你多关照。}
\end{dialogue}

\begin{culture}
  初次见面时，\jp{「よろしくお願いします」}不是逐字对应某一句中文，
  而是表达希望今后关系顺利、请对方关照的固定寒暄语。对方常用
  \jp{「こちらこそ」}回应，意思接近“我才要请您多关照”。
\end{culture}

\subsection*{练习与复盘}

\begin{practice}
  将下列句子翻译成日语。
  \begin{enumerate}
    \item 我是中国人。
    \item 铃木先生是老师吗？
    \item 我不是大学生。
  \end{enumerate}
  \answerline[3]
\end{practice}

\begin{practice}
  请用本课句型写一段三到五句的自我介绍。
  \answerline[5]
\end{practice}